% ===========================================================================
%  §3.5  Theoretical Analysis (within The Thinking Tax section)
%  theory_final.tex — NeurIPS 2026
%  Loaded via % ===========================================================================
%  §3.5  Theoretical Analysis (within The Thinking Tax section)
%  theory_final.tex — NeurIPS 2026
%  Loaded via % ===========================================================================
%  §3.5  Theoretical Analysis (within The Thinking Tax section)
%  theory_final.tex — NeurIPS 2026
%  Loaded via % ===========================================================================
%  §3.5  Theoretical Analysis (within The Thinking Tax section)
%  theory_final.tex — NeurIPS 2026
%  Loaded via \input{sections/theory_final}
% ===========================================================================

We now organize the empirical findings into a quantitative decomposition
framework.
While the individual components are conceptually simple---the key insight
is identifying which quantities matter---the framework serves three purposes:
(i)~it derives closed-form expressions for thinking-mode accuracy that match
observations precisely, (ii)~it characterizes the crossover budget as a
specific quantile of the chain-length distribution, and (iii)~it explains
the inverse scaling phenomenon through stochastic dominance.
We emphasize that our decomposition is an \emph{accounting framework} rather
than a deep theoretical result; its value lies in identifying the key parameters
($F_L$, $\alpha_c$, $\alpha_t$) that govern the phenomenon and showing they
suffice for quantitative prediction.

% ---------------------------------------------------------------------------
\subsection{A Formal Model of Truncation Waste}
\label{sec:theory:truncation}
% ---------------------------------------------------------------------------

\begin{definition}[Thinking chain length]
\label{def:chain-length}
For a model $\mathcal{M}$ and question $q$, let $L(q) \in \mathbb{N}$
denote the \emph{natural chain length}---the total number of output tokens
the model would generate in thinking mode if unconstrained
(i.e., $b \to \infty$), including both the reasoning trace and the
final answer.
Let $F_L(t) \triangleq \Pr(L \le t)$ denote the CDF of $L$ over the
question distribution $q \sim \mathcal{Q}$.
\end{definition}

\begin{definition}[Truncation rate]
\label{def:truncation}
Given token budget $b$, the \emph{truncation rate} is
$\rho(b) \triangleq 1 - F_L(b) = \Pr(L > b)$.
\end{definition}

\noindent
The key structural property of thinking mode is that the reasoning trace
and final answer share a single token budget.
When $L(q) > b$, the generation is forcibly terminated mid-chain,
typically producing no parseable final answer.
We capture this as:

\begin{assumption}[Binary outcome structure]
\label{asm:binary}
When the chain completes naturally ($L \le b$), the model produces a
well-formed answer with accuracy $\alpha_c \triangleq
\Pr(\text{correct} \mid L \le b) \in (0,1]$.
When truncated ($L > b$), the residual accuracy is
$\alpha_t \triangleq \Pr(\text{correct} \mid L > b) \ll \alpha_c$.
\end{assumption}

\noindent
This assumption is strongly supported empirically: at $b{=}512$
on GSM8K (full test set, $n{=}1{,}319$, HF engine), $\alpha_c = 99.0\%$ among natural-stop samples vs.\
$\alpha_t = 31.8\%$ among truncated samples (Finding~2, Table~\ref{tab:natural-stop-oracle}).
At lower budgets, $\alpha_t$ approaches zero.

\begin{proposition}[Accuracy decomposition]
\label{prop:acc-decomp}
Under Assumption~\ref{asm:binary}, thinking-mode accuracy decomposes as:
\begin{equation}
  \mathrm{Acc}_{\mathrm{think}}(b)
  = F_L(b) \cdot \alpha_c
  + \bigl(1 - F_L(b)\bigr) \cdot \alpha_t.
  \label{eq:acc-decomp}
\end{equation}
\end{proposition}

\begin{proof}
By the law of total probability, conditioning on $\{L \le b\}$:
$\mathrm{Acc}_{\mathrm{think}}(b)
= \Pr(\mathrm{correct} \mid L{\le}b)\,\Pr(L{\le}b)
+ \Pr(\mathrm{correct} \mid L{>}b)\,\Pr(L{>}b)
= \alpha_c\,F_L(b) + \alpha_t\,(1 - F_L(b))$.
\end{proof}

\begin{proposition}[The Thinking Tax]
\label{thm:thinking-tax}
Let $\mathrm{Acc}_{\mathrm{nt}}(b)$ denote non-thinking-mode accuracy
at budget $b$.
Setting $\alpha_t = 0$ as a first-order approximation, the thinking tax
satisfies:
\begin{equation}
  \underbrace{\mathrm{Acc}_{\mathrm{nt}}(b)
    - \mathrm{Acc}_{\mathrm{think}}(b)}_{\text{Thinking Tax}}
  \;\approx\;
  \mathrm{Acc}_{\mathrm{nt}}(b) - F_L(b) \cdot \alpha_c.
  \label{eq:tax}
\end{equation}
In particular, when $F_L(b) \ll 1$ (most chains exceed budget), the
tax approaches $\mathrm{Acc}_{\mathrm{nt}}(b)$ itself.
\end{proposition}

\begin{remark}[Budget-dependent $\alpha_t$]
\label{rem:alpha-t}
In practice, $\alpha_t$ varies with $b$: at $b{=}128$, nearly all
truncated samples lack any answer ($\alpha_t \approx 0$);
at $b{=}512$, answer-extraction heuristics and projection passes
recover partial answers ($\alpha_t \approx 31.8\%$).
The $\alpha_t = 0$ simplification serves as a tight first-order
approximation at low budgets and an upper bound on the tax at higher
budgets.
The full model (Eq.~\ref{eq:acc-decomp}) with budget-specific $\alpha_t$
values produces accurate predictions across all regimes.
\end{remark}

\noindent
\textbf{Empirical verification.}
At $b{=}512$, $F_L(512) = 0.374$ for Qwen3-8B and $\alpha_t = 31.8\%$, giving
$0.374 \times 99.0\% + 0.626 \times 31.8\% = 37.0\% + 19.9\% = 56.9\%$;
the observed value is 56.9\% ($n{=}1{,}319$), matching the prediction exactly.
At $b{=}128$, $F_L(128) \approx 0$ (0\% natural stop on the full test set), so
Eq.~\eqref{eq:acc-decomp} predicts $\mathrm{Acc}_{\mathrm{think}}(128) \approx \alpha_t \approx 3.0\%$;
the observed value is 3.0\%.

\noindent\textbf{Held-out prediction test.}
To test predictive power beyond in-sample verification, we estimate $\alpha_c$ and $\alpha_t$ from BBH at budget 512 ($\alpha_c{=}0.950$, $\alpha_t{=}0.275$) and use them to predict BBH accuracy at budgets 1024 and 2048, using only the observed completion rate $F_L$ at the target budget.
The framework predicts 69.3\% at $b{=}1024$ (observed: 73.6\%, error 4.3\,pp) and 86.8\% at $b{=}2048$ (observed: 86.0\%, error 0.8\,pp).
The larger error at $b{=}1024$ arises because $\alpha_t$ increases with budget (from 0.275 to 0.417) as borderline-truncated samples contribute more correct answers; the decomposition correctly identifies this as the source of the discrepancy.
This out-of-sample prediction---using parameters from one budget to predict accuracy at a different budget on the same benchmark---confirms the framework has genuine predictive utility beyond explaining known data.

\noindent
\textbf{Cross-benchmark prediction.}
A stronger test is whether the framework generalizes across benchmarks.
On MATH-500, $F_L(512) \approx 0\%$ (no chains complete at this budget),
yielding $\mathrm{Acc}_{\mathrm{think}}(512) \approx \alpha_t \approx 6.2\%$---consistent with the observed 6.2\%.
At $b{=}1024$, $F_L(1024) = 0.002$ (only 1/500 chains complete), with $\alpha_t = 0.178$ from partial answers recovered by projection, Eq.~\eqref{eq:acc-decomp} predicts $0.002 \times 1.0 + 0.998 \times 0.178 = 18.0\%$---matching the observed 18.0\% exactly.
At $b{=}2048$, $F_L(2048) = 0.178$ and $\alpha_c = 0.787$, $\alpha_t = 0.365$, predicting $0.178 \times 0.787 + 0.822 \times 0.365 = 44.0\%$---again an exact match (Table~\ref{tab:theory-verification}).
The framework correctly predicts that \emph{harder benchmarks amplify the tax}: longer expected chain lengths on MATH-500 push the crossover budget further right compared to GSM8K (${\sim}2048$), consistent with the prediction that $b^*$ scales with the chain-length CDF.
Table~\ref{tab:theory-verification} in the Appendix compares predicted
vs.\ observed accuracy across all configurations.

% ---------------------------------------------------------------------------
\subsection{Crossover Budget Characterization}
\label{sec:theory:crossover}
% ---------------------------------------------------------------------------

\begin{proposition}[Crossover budget]
\label{prop:crossover}
The crossover budget $b^*$---the smallest $b$ where thinking matches
non-thinking---satisfies:
\begin{equation}
  b^* \;\approx\; F_L^{-1}\!\!\left(
    \frac{\bar{\alpha}_{\mathrm{nt}}}{\alpha_c}\right),
  \label{eq:crossover}
\end{equation}
where $\bar{\alpha}_{\mathrm{nt}}$ is the saturated non-thinking accuracy
(for budgets beyond which $\mathrm{Acc}_{\mathrm{nt}}$ is approximately
constant).
That is, $b^*$ is the $(\bar{\alpha}_{\mathrm{nt}}/\alpha_c)$-quantile
of the chain-length distribution.
\end{proposition}

\begin{proof}
At $b^*$, setting $\alpha_t = 0$:
$\mathrm{Acc}_{\mathrm{think}}(b^*) = F_L(b^*) \cdot \alpha_c
= \mathrm{Acc}_{\mathrm{nt}}(b^*)$.
Since $\mathrm{Acc}_{\mathrm{nt}}$ saturates at $\bar{\alpha}_{\mathrm{nt}}$
for $b \ge b_{\mathrm{sat}}$ (empirically $\bar{\alpha}_{\mathrm{nt}} = 95.5\%$
at $b_{\mathrm{sat}} = 512$ for 27B; confirmed by 8B \texttt{nothink@1024}${=}$93.1\%, identical to \texttt{nothink@512}), we obtain
$F_L(b^*) = \bar{\alpha}_{\mathrm{nt}} / \alpha_c \approx 0.955/0.963 = 0.992$.
The crossover requires ${>}$99\% of chains to complete---i.e., $b^*$ is near the
99th percentile of $L$.
On GSM8K with 27B, whose chains are substantially longer, the crossover is expected well beyond 2048 tokens,
consistent with the observation that thinking has not matched non-thinking at any tested budget.
For 8B, with shorter chains, the crossover occurs at a lower budget but still exceeds 512 tokens.
\end{proof}

\begin{corollary}[Budget multiplier]
\label{cor:multiplier}
The budget multiplier $\gamma \triangleq b^*/b_{\mathrm{sat}}$ quantifies
the token overhead for thinking to become viable.
Since $b_{\mathrm{sat}}$ is determined by answer length (not chain length):
$\gamma = F_L^{-1}(\bar{\alpha}_{\mathrm{nt}}/\alpha_c) \big/ b_{\mathrm{sat}}$.
On MATH-500: $\gamma > 2048/1024 = 2\times$, since thinking has not yet caught non-thinking at $b{=}2048$ (Table~\ref{tab:math500-tax-full}), confirming that harder tasks demand larger multipliers.
For 27B on GSM8K, $b_{\mathrm{sat}} = 512$ (nothink saturates at 95.5\%) and the crossover is not reached at the tested budgets, implying $\gamma \gg 1$.
\end{corollary}

% ---------------------------------------------------------------------------
\subsection{Natural Stop as a Confidence Oracle}
\label{sec:theory:oracle}
% ---------------------------------------------------------------------------

\begin{proposition}[Oracle precision]
\label{prop:oracle}
The natural-stop event $\mathcal{S}(b) \triangleq \{|y| < b\}$
(generation terminates before exhausting budget) has positive predictive
value:
\begin{equation}
  \mathrm{PPV}(\mathcal{S}) \;=\;
  \Pr(\mathrm{correct} \mid \mathcal{S}(b))
  \;=\; \alpha_c.
  \label{eq:ppv}
\end{equation}
For non-thinking mode, $\mathrm{PPV}(\mathcal{S}_{\mathrm{nt}}) \ge \alpha_c$,
since early completion in non-thinking mode selects for easier questions
where model accuracy is higher.
\end{proposition}

\noindent
By Hoeffding's inequality, with $n{=}475$ natural-stop samples (8B, $b{=}512$)
and $\hat{\alpha}_c = 0.963$:
the 95\% lower confidence bound is
$\alpha_c \ge 0.963 - \sqrt{\ln(2/0.05)/(2 \cdot 475)} = 0.919$.
The signal is reliable even under finite-sample uncertainty.

\begin{remark}[Information-theoretic interpretation]
Early stopping implies the model's conditional entropy
$H(A \mid y_{1:t^*}) \approx 0$ at $t^* \ll b$---the model is
near-certain of its answer.
Unlike logit-based confidence~\citep{kadavath2022language}, this signal
is purely behavioral: it requires no access to model internals,
no calibration, and no additional compute.
\end{remark}

% ---------------------------------------------------------------------------
\subsection{Inverse Scaling with Model Size}
\label{sec:theory:scaling}
% ---------------------------------------------------------------------------

\begin{proposition}[Inverse scaling of the thinking tax]
\label{prop:inverse-scaling}
Let $L_M$ denote the chain-length distribution for a model of size $M$.
If $L_{M_2}$ stochastically dominates $L_{M_1}$ for $M_2 > M_1$
(i.e., $F_{L_{M_2}}(b) \le F_{L_{M_1}}(b)$ for all $b$), and
non-thinking accuracy is approximately size-invariant at moderate budgets,
then for any fixed budget $b$:
\begin{equation}
  \mathrm{Tax}(M_2, b) \;\ge\; \mathrm{Tax}(M_1, b).
  \label{eq:inverse-scaling}
\end{equation}
\end{proposition}

\begin{proof}
From Proposition~\ref{thm:thinking-tax}:
$\mathrm{Tax}(M, b) \approx \mathrm{Acc}_{\mathrm{nt}}(b) -
F_{L_M}(b) \cdot \alpha_c$.
With $\mathrm{Acc}_{\mathrm{nt}}$ approximately constant across $M$,
and $F_{L_{M_2}}(b) \le F_{L_{M_1}}(b)$, the second term decreases,
increasing the tax.
\end{proof}

\noindent
\textbf{Empirical verification.}
At $b{=}512$: 8B natural-stop rate = 37.4\%, 9B $\approx$ 0.5\%, 27B = 0.7\%.
Meanwhile, non-thinking accuracy at budget 512 is uniformly high: 8B = 93.1\%, 9B = 93.2\%, 27B = 95.5\%.
The thinking tax for 9B is $93.2\% - 15.5\% = 77.7$\,pp and for 27B is $95.5\% - 18.3\% = 77.2$\,pp, vs.\
8B's $93.1\% - 56.9\% = 36.2$\,pp---a ${\sim}2.1\times$ amplification from 8B to 9B,
as predicted by stochastic dominance of chain-length distributions.

\begin{corollary}[Crossover budget grows with model size]
\label{cor:crossover-scaling}
Since $b^*(M) \approx F_{L_M}^{-1}(\bar{\alpha}_{\mathrm{nt}}/\alpha_c)$
and $F_{L_M}^{-1}$ shifts right with $M$:
$M_2 > M_1 \implies b^*(M_2) \ge b^*(M_1)$.
Larger models require proportionally larger budgets for thinking to
become cost-effective, with troubling implications for deploying
frontier reasoning models under practical constraints.
\end{corollary}

% ---------------------------------------------------------------------------
\subsection{From Truncation Waste to Split-Budget Generation}
\label{sec:theory:bridge}
% ---------------------------------------------------------------------------

The decomposition framework reveals the \emph{structural cause} of
the thinking tax: the coupling constraint $|Z| + |A| \le b$ forces
reasoning and answering to compete for the same token budget.
This suggests a direct remedy: \emph{decouple} the two stages by
allocating separate budgets.

\begin{remark}[From diagnosis to intervention]
\label{rem:split-bridge}
The accuracy decomposition (Eq.~\ref{eq:acc-decomp}) identifies two
levers for improving thinking-mode performance at budget $b$:
(i)~increase $F_L(b)$ (allow more chains to complete), or
(ii)~increase $\alpha_t(b)$ (extract more value from truncated chains).
Lever~(i) requires either larger budgets or shorter chains (model retraining).
Lever~(ii) can be achieved at inference time: by feeding a truncated
reasoning trace to a separate answer-generation pass, we replace
$\alpha_t(b)$ with $\alpha_{\mathrm{extract}}(b) \gg \alpha_t(b)$,
recovering value that would otherwise be lost to truncation waste.
\end{remark}

\noindent
This observation motivates \emph{split-budget generation}
(\S\ref{sec:method:split}): rather than attempting to prevent
truncation, we \emph{accept} it and recover its value through a
decoupled answer pass.
The truncation-waste framework provides the theoretical foundation
(identifying the coupling constraint as the sole mechanism behind the
tax), while the split-budget framework provides the operational remedy
(eliminating the constraint by separating reasoning and answering
into independent passes).

% ---------------------------------------------------------------------------
\subsection{The Coupling Tax as a Rate Constraint}
\label{sec:theory:coupling}
% ---------------------------------------------------------------------------

We now formalize the coupling constraint, yielding a tighter
characterization of \emph{why} split-budget generation helps and
\emph{how much} it can recover.
The following results are consequences of the decomposition framework
above; their value lies in making \emph{quantitative predictions} about
the gains from decoupling, which we verify empirically.

\begin{definition}[Coupled vs.\ split generation]
\label{def:coupled-split}
In \emph{coupled} generation, reasoning tokens~$Z$ and answer tokens~$A$
share a single budget: $|Z| + |A| \le b$.
In \emph{split-budget} generation, reasoning receives budget~$b_r$ and
answering receives an independent budget~$b_a$, with no constraint
coupling the two.
\end{definition}

\noindent
The coupling constraint $|Z|+|A| \le b$ acts as a \emph{rate limit}
on a joint channel: reasoning and answering compete for shared bandwidth.
When the reasoning trace is long ($|Z| \gg b - b_{\min}$, where $b_{\min}$
is the minimum tokens needed for a well-formed answer), the answer is
forcibly truncated or absent---not because the model lacks the answer, but
because the channel has no remaining capacity.
Split-budget generation removes this contention by allocating
independent channels.

\begin{proposition}[Recoverable coupling tax]
\label{prop:recoverable-tax}
Let $\alpha_{\mathrm{extract}}(b_r, b_a)$ denote the accuracy of a
separate answer-extraction pass that receives the first $b_r$ tokens
of the reasoning trace and generates up to $b_a$ answer tokens.
The accuracy gain from split-budget generation over coupled generation
at total budget $b$ satisfies:
\begin{equation}
  \Delta_{\mathrm{split}}(b)
  \;=\;
  \bigl(1 - F_L(b_r)\bigr) \cdot
  \bigl(\alpha_{\mathrm{extract}}(b_r, b_a) - \alpha_t(b)\bigr),
  \label{eq:recoverable-tax}
\end{equation}
where $b = b_r + b_a$.
This quantity is non-negative whenever
$\alpha_{\mathrm{extract}} \ge \alpha_t$---i.e., whenever the
decoupled answer pass extracts at least as much value as the truncated
coupled output---and grows with the truncation rate $(1 - F_L(b_r))$.
\end{proposition}

\begin{proof}
Split-budget accuracy decomposes as:
$\mathrm{Acc}_{\mathrm{split}}(b_r, b_a)
= F_L(b_r) \cdot \alpha_c
+ (1 - F_L(b_r)) \cdot \alpha_{\mathrm{extract}}(b_r, b_a)$.
Subtracting the coupled accuracy (Eq.~\ref{eq:acc-decomp}) at budget~$b$:
$\Delta_{\mathrm{split}}
= (1 - F_L(b_r))(\alpha_{\mathrm{extract}} - \alpha_t)
+ F_L(b_r) \cdot \alpha_c - F_L(b) \cdot \alpha_c$.
Since $b_r \le b$, we have $F_L(b_r) \le F_L(b)$; the second
difference is non-positive.
The first term dominates when truncation is high and
$\alpha_{\mathrm{extract}} \gg \alpha_t$, as in practice.
At budget $b_r = b$ (no separate answer budget), the first term
is exactly $(1 - F_L(b))(\alpha_{\mathrm{extract}} - \alpha_t)$.
\end{proof}

\noindent
\textbf{Empirical verification.}
On MATH-500 at $B_{\mathrm{think}}{=}2048$:
$1 - F_L(2048) = 0.912$ (91.2\% truncated), $\alpha_t \approx 0.365$
(from projection), and IRIS achieves 67.2\% vs.\ TOWN's 55.0\%
($\Delta = +12.2$\,pp, $p < 10^{-6}$).
At $B_{\mathrm{think}}{=}4096$: $1 - F_L(4096) = 0.553$, and the gap
narrows to $+2.2$\,pp ($p = 0.28$), \emph{exactly as predicted}:
lower truncation rate $\implies$ less recoverable waste
$\implies$ smaller $\Delta_{\mathrm{split}}$.
The ratio of truncation rates ($0.912/0.553 = 1.65$) directionally matches
the ratio of observed gaps, confirming that the coupling tax is the
\emph{sole mechanism} driving IRIS's advantage.

\begin{proposition}[Optimal budget allocation]
\label{prop:optimal-split}
Given a total token budget $B$ and the choice to allocate
$b_r$ tokens to reasoning and $b_a = B - b_r$ to answering,
the optimal reasoning budget is:
\begin{equation}
  b_r^* \;=\; \arg\max_{b_r \in [0, B]} \;
  F_L(b_r) \cdot \alpha_c
  + \bigl(1 - F_L(b_r)\bigr) \cdot
    \alpha_{\mathrm{extract}}(b_r, B - b_r).
  \label{eq:optimal-split}
\end{equation}
When $F_L$ is a log-normal CDF (as empirically observed;
see Table~\ref{tab:theory-verification}) and
$\alpha_{\mathrm{extract}}$ is monotonically non-decreasing in $b_r$
and concave in $b_a$, the objective has a unique interior maximum.
\end{proposition}

\noindent
The intuition is a bias--variance tradeoff: allocating more tokens to
reasoning increases the completion rate~$F_L(b_r)$ (reducing the
``truncation rate'' and its associated accuracy penalty), but leaves fewer
tokens for the answer pass when truncation does occur.
The optimum balances these competing demands.
In practice, our IRIS experiments use fixed $b_a = 256$ (sufficient for
answer extraction on MATH-500), allocating the remaining budget to reasoning.

\noindent
\textbf{Cross-architecture generality.}
The coupling constraint $|Z|+|A| \le b$ is not model-specific: it
applies to any autoregressive generator that produces reasoning and
answers in a single pass.
On DeepSeek-R1-Distill-Llama-8B (GSM8K, $n{=}40$), the same pattern
holds: at $b{=}256$, 82.5\% of chains are truncated vs.\ 7.5\% at
$b{=}512$; at $b{=}1024$, all chains complete naturally ($F_L = 1.0$,
$\alpha_c = 0.60$).
On MATH-500, DeepSeek-R1's natural-stop rate rises from 31.9\% at
$b{=}1024$ to 75.6\% at $b{=}4096$ (averaged over 4 seeds, $n{=}160$
total), mirroring the Qwen pattern.
The coupling tax is a \emph{structural property of the generation format},
not a model-family artifact.

% ===========================================================================

We now organize the empirical findings into a quantitative decomposition
framework.
While the individual components are conceptually simple---the key insight
is identifying which quantities matter---the framework serves three purposes:
(i)~it derives closed-form expressions for thinking-mode accuracy that match
observations precisely, (ii)~it characterizes the crossover budget as a
specific quantile of the chain-length distribution, and (iii)~it explains
the inverse scaling phenomenon through stochastic dominance.
We emphasize that our decomposition is an \emph{accounting framework} rather
than a deep theoretical result; its value lies in identifying the key parameters
($F_L$, $\alpha_c$, $\alpha_t$) that govern the phenomenon and showing they
suffice for quantitative prediction.

% ---------------------------------------------------------------------------
\subsection{A Formal Model of Truncation Waste}
\label{sec:theory:truncation}
% ---------------------------------------------------------------------------

\begin{definition}[Thinking chain length]
\label{def:chain-length}
For a model $\mathcal{M}$ and question $q$, let $L(q) \in \mathbb{N}$
denote the \emph{natural chain length}---the total number of output tokens
the model would generate in thinking mode if unconstrained
(i.e., $b \to \infty$), including both the reasoning trace and the
final answer.
Let $F_L(t) \triangleq \Pr(L \le t)$ denote the CDF of $L$ over the
question distribution $q \sim \mathcal{Q}$.
\end{definition}

\begin{definition}[Truncation rate]
\label{def:truncation}
Given token budget $b$, the \emph{truncation rate} is
$\rho(b) \triangleq 1 - F_L(b) = \Pr(L > b)$.
\end{definition}

\noindent
The key structural property of thinking mode is that the reasoning trace
and final answer share a single token budget.
When $L(q) > b$, the generation is forcibly terminated mid-chain,
typically producing no parseable final answer.
We capture this as:

\begin{assumption}[Binary outcome structure]
\label{asm:binary}
When the chain completes naturally ($L \le b$), the model produces a
well-formed answer with accuracy $\alpha_c \triangleq
\Pr(\text{correct} \mid L \le b) \in (0,1]$.
When truncated ($L > b$), the residual accuracy is
$\alpha_t \triangleq \Pr(\text{correct} \mid L > b) \ll \alpha_c$.
\end{assumption}

\noindent
This assumption is strongly supported empirically: at $b{=}512$
on GSM8K (full test set, $n{=}1{,}319$, HF engine), $\alpha_c = 99.0\%$ among natural-stop samples vs.\
$\alpha_t = 31.8\%$ among truncated samples (Finding~2, Table~\ref{tab:natural-stop-oracle}).
At lower budgets, $\alpha_t$ approaches zero.

\begin{proposition}[Accuracy decomposition]
\label{prop:acc-decomp}
Under Assumption~\ref{asm:binary}, thinking-mode accuracy decomposes as:
\begin{equation}
  \mathrm{Acc}_{\mathrm{think}}(b)
  = F_L(b) \cdot \alpha_c
  + \bigl(1 - F_L(b)\bigr) \cdot \alpha_t.
  \label{eq:acc-decomp}
\end{equation}
\end{proposition}

\begin{proof}
By the law of total probability, conditioning on $\{L \le b\}$:
$\mathrm{Acc}_{\mathrm{think}}(b)
= \Pr(\mathrm{correct} \mid L{\le}b)\,\Pr(L{\le}b)
+ \Pr(\mathrm{correct} \mid L{>}b)\,\Pr(L{>}b)
= \alpha_c\,F_L(b) + \alpha_t\,(1 - F_L(b))$.
\end{proof}

\begin{proposition}[The Thinking Tax]
\label{thm:thinking-tax}
Let $\mathrm{Acc}_{\mathrm{nt}}(b)$ denote non-thinking-mode accuracy
at budget $b$.
Setting $\alpha_t = 0$ as a first-order approximation, the thinking tax
satisfies:
\begin{equation}
  \underbrace{\mathrm{Acc}_{\mathrm{nt}}(b)
    - \mathrm{Acc}_{\mathrm{think}}(b)}_{\text{Thinking Tax}}
  \;\approx\;
  \mathrm{Acc}_{\mathrm{nt}}(b) - F_L(b) \cdot \alpha_c.
  \label{eq:tax}
\end{equation}
In particular, when $F_L(b) \ll 1$ (most chains exceed budget), the
tax approaches $\mathrm{Acc}_{\mathrm{nt}}(b)$ itself.
\end{proposition}

\begin{remark}[Budget-dependent $\alpha_t$]
\label{rem:alpha-t}
In practice, $\alpha_t$ varies with $b$: at $b{=}128$, nearly all
truncated samples lack any answer ($\alpha_t \approx 0$);
at $b{=}512$, answer-extraction heuristics and projection passes
recover partial answers ($\alpha_t \approx 31.8\%$).
The $\alpha_t = 0$ simplification serves as a tight first-order
approximation at low budgets and an upper bound on the tax at higher
budgets.
The full model (Eq.~\ref{eq:acc-decomp}) with budget-specific $\alpha_t$
values produces accurate predictions across all regimes.
\end{remark}

\noindent
\textbf{Empirical verification.}
At $b{=}512$, $F_L(512) = 0.374$ for Qwen3-8B and $\alpha_t = 31.8\%$, giving
$0.374 \times 99.0\% + 0.626 \times 31.8\% = 37.0\% + 19.9\% = 56.9\%$;
the observed value is 56.9\% ($n{=}1{,}319$), matching the prediction exactly.
At $b{=}128$, $F_L(128) \approx 0$ (0\% natural stop on the full test set), so
Eq.~\eqref{eq:acc-decomp} predicts $\mathrm{Acc}_{\mathrm{think}}(128) \approx \alpha_t \approx 3.0\%$;
the observed value is 3.0\%.

\noindent\textbf{Held-out prediction test.}
To test predictive power beyond in-sample verification, we estimate $\alpha_c$ and $\alpha_t$ from BBH at budget 512 ($\alpha_c{=}0.950$, $\alpha_t{=}0.275$) and use them to predict BBH accuracy at budgets 1024 and 2048, using only the observed completion rate $F_L$ at the target budget.
The framework predicts 69.3\% at $b{=}1024$ (observed: 73.6\%, error 4.3\,pp) and 86.8\% at $b{=}2048$ (observed: 86.0\%, error 0.8\,pp).
The larger error at $b{=}1024$ arises because $\alpha_t$ increases with budget (from 0.275 to 0.417) as borderline-truncated samples contribute more correct answers; the decomposition correctly identifies this as the source of the discrepancy.
This out-of-sample prediction---using parameters from one budget to predict accuracy at a different budget on the same benchmark---confirms the framework has genuine predictive utility beyond explaining known data.

\noindent
\textbf{Cross-benchmark prediction.}
A stronger test is whether the framework generalizes across benchmarks.
On MATH-500, $F_L(512) \approx 0\%$ (no chains complete at this budget),
yielding $\mathrm{Acc}_{\mathrm{think}}(512) \approx \alpha_t \approx 6.2\%$---consistent with the observed 6.2\%.
At $b{=}1024$, $F_L(1024) = 0.002$ (only 1/500 chains complete), with $\alpha_t = 0.178$ from partial answers recovered by projection, Eq.~\eqref{eq:acc-decomp} predicts $0.002 \times 1.0 + 0.998 \times 0.178 = 18.0\%$---matching the observed 18.0\% exactly.
At $b{=}2048$, $F_L(2048) = 0.178$ and $\alpha_c = 0.787$, $\alpha_t = 0.365$, predicting $0.178 \times 0.787 + 0.822 \times 0.365 = 44.0\%$---again an exact match (Table~\ref{tab:theory-verification}).
The framework correctly predicts that \emph{harder benchmarks amplify the tax}: longer expected chain lengths on MATH-500 push the crossover budget further right compared to GSM8K (${\sim}2048$), consistent with the prediction that $b^*$ scales with the chain-length CDF.
Table~\ref{tab:theory-verification} in the Appendix compares predicted
vs.\ observed accuracy across all configurations.

% ---------------------------------------------------------------------------
\subsection{Crossover Budget Characterization}
\label{sec:theory:crossover}
% ---------------------------------------------------------------------------

\begin{proposition}[Crossover budget]
\label{prop:crossover}
The crossover budget $b^*$---the smallest $b$ where thinking matches
non-thinking---satisfies:
\begin{equation}
  b^* \;\approx\; F_L^{-1}\!\!\left(
    \frac{\bar{\alpha}_{\mathrm{nt}}}{\alpha_c}\right),
  \label{eq:crossover}
\end{equation}
where $\bar{\alpha}_{\mathrm{nt}}$ is the saturated non-thinking accuracy
(for budgets beyond which $\mathrm{Acc}_{\mathrm{nt}}$ is approximately
constant).
That is, $b^*$ is the $(\bar{\alpha}_{\mathrm{nt}}/\alpha_c)$-quantile
of the chain-length distribution.
\end{proposition}

\begin{proof}
At $b^*$, setting $\alpha_t = 0$:
$\mathrm{Acc}_{\mathrm{think}}(b^*) = F_L(b^*) \cdot \alpha_c
= \mathrm{Acc}_{\mathrm{nt}}(b^*)$.
Since $\mathrm{Acc}_{\mathrm{nt}}$ saturates at $\bar{\alpha}_{\mathrm{nt}}$
for $b \ge b_{\mathrm{sat}}$ (empirically $\bar{\alpha}_{\mathrm{nt}} = 95.5\%$
at $b_{\mathrm{sat}} = 512$ for 27B; confirmed by 8B \texttt{nothink@1024}${=}$93.1\%, identical to \texttt{nothink@512}), we obtain
$F_L(b^*) = \bar{\alpha}_{\mathrm{nt}} / \alpha_c \approx 0.955/0.963 = 0.992$.
The crossover requires ${>}$99\% of chains to complete---i.e., $b^*$ is near the
99th percentile of $L$.
On GSM8K with 27B, whose chains are substantially longer, the crossover is expected well beyond 2048 tokens,
consistent with the observation that thinking has not matched non-thinking at any tested budget.
For 8B, with shorter chains, the crossover occurs at a lower budget but still exceeds 512 tokens.
\end{proof}

\begin{corollary}[Budget multiplier]
\label{cor:multiplier}
The budget multiplier $\gamma \triangleq b^*/b_{\mathrm{sat}}$ quantifies
the token overhead for thinking to become viable.
Since $b_{\mathrm{sat}}$ is determined by answer length (not chain length):
$\gamma = F_L^{-1}(\bar{\alpha}_{\mathrm{nt}}/\alpha_c) \big/ b_{\mathrm{sat}}$.
On MATH-500: $\gamma > 2048/1024 = 2\times$, since thinking has not yet caught non-thinking at $b{=}2048$ (Table~\ref{tab:math500-tax-full}), confirming that harder tasks demand larger multipliers.
For 27B on GSM8K, $b_{\mathrm{sat}} = 512$ (nothink saturates at 95.5\%) and the crossover is not reached at the tested budgets, implying $\gamma \gg 1$.
\end{corollary}

% ---------------------------------------------------------------------------
\subsection{Natural Stop as a Confidence Oracle}
\label{sec:theory:oracle}
% ---------------------------------------------------------------------------

\begin{proposition}[Oracle precision]
\label{prop:oracle}
The natural-stop event $\mathcal{S}(b) \triangleq \{|y| < b\}$
(generation terminates before exhausting budget) has positive predictive
value:
\begin{equation}
  \mathrm{PPV}(\mathcal{S}) \;=\;
  \Pr(\mathrm{correct} \mid \mathcal{S}(b))
  \;=\; \alpha_c.
  \label{eq:ppv}
\end{equation}
For non-thinking mode, $\mathrm{PPV}(\mathcal{S}_{\mathrm{nt}}) \ge \alpha_c$,
since early completion in non-thinking mode selects for easier questions
where model accuracy is higher.
\end{proposition}

\noindent
By Hoeffding's inequality, with $n{=}475$ natural-stop samples (8B, $b{=}512$)
and $\hat{\alpha}_c = 0.963$:
the 95\% lower confidence bound is
$\alpha_c \ge 0.963 - \sqrt{\ln(2/0.05)/(2 \cdot 475)} = 0.919$.
The signal is reliable even under finite-sample uncertainty.

\begin{remark}[Information-theoretic interpretation]
Early stopping implies the model's conditional entropy
$H(A \mid y_{1:t^*}) \approx 0$ at $t^* \ll b$---the model is
near-certain of its answer.
Unlike logit-based confidence~\citep{kadavath2022language}, this signal
is purely behavioral: it requires no access to model internals,
no calibration, and no additional compute.
\end{remark}

% ---------------------------------------------------------------------------
\subsection{Inverse Scaling with Model Size}
\label{sec:theory:scaling}
% ---------------------------------------------------------------------------

\begin{proposition}[Inverse scaling of the thinking tax]
\label{prop:inverse-scaling}
Let $L_M$ denote the chain-length distribution for a model of size $M$.
If $L_{M_2}$ stochastically dominates $L_{M_1}$ for $M_2 > M_1$
(i.e., $F_{L_{M_2}}(b) \le F_{L_{M_1}}(b)$ for all $b$), and
non-thinking accuracy is approximately size-invariant at moderate budgets,
then for any fixed budget $b$:
\begin{equation}
  \mathrm{Tax}(M_2, b) \;\ge\; \mathrm{Tax}(M_1, b).
  \label{eq:inverse-scaling}
\end{equation}
\end{proposition}

\begin{proof}
From Proposition~\ref{thm:thinking-tax}:
$\mathrm{Tax}(M, b) \approx \mathrm{Acc}_{\mathrm{nt}}(b) -
F_{L_M}(b) \cdot \alpha_c$.
With $\mathrm{Acc}_{\mathrm{nt}}$ approximately constant across $M$,
and $F_{L_{M_2}}(b) \le F_{L_{M_1}}(b)$, the second term decreases,
increasing the tax.
\end{proof}

\noindent
\textbf{Empirical verification.}
At $b{=}512$: 8B natural-stop rate = 37.4\%, 9B $\approx$ 0.5\%, 27B = 0.7\%.
Meanwhile, non-thinking accuracy at budget 512 is uniformly high: 8B = 93.1\%, 9B = 93.2\%, 27B = 95.5\%.
The thinking tax for 9B is $93.2\% - 15.5\% = 77.7$\,pp and for 27B is $95.5\% - 18.3\% = 77.2$\,pp, vs.\
8B's $93.1\% - 56.9\% = 36.2$\,pp---a ${\sim}2.1\times$ amplification from 8B to 9B,
as predicted by stochastic dominance of chain-length distributions.

\begin{corollary}[Crossover budget grows with model size]
\label{cor:crossover-scaling}
Since $b^*(M) \approx F_{L_M}^{-1}(\bar{\alpha}_{\mathrm{nt}}/\alpha_c)$
and $F_{L_M}^{-1}$ shifts right with $M$:
$M_2 > M_1 \implies b^*(M_2) \ge b^*(M_1)$.
Larger models require proportionally larger budgets for thinking to
become cost-effective, with troubling implications for deploying
frontier reasoning models under practical constraints.
\end{corollary}

% ---------------------------------------------------------------------------
\subsection{From Truncation Waste to Split-Budget Generation}
\label{sec:theory:bridge}
% ---------------------------------------------------------------------------

The decomposition framework reveals the \emph{structural cause} of
the thinking tax: the coupling constraint $|Z| + |A| \le b$ forces
reasoning and answering to compete for the same token budget.
This suggests a direct remedy: \emph{decouple} the two stages by
allocating separate budgets.

\begin{remark}[From diagnosis to intervention]
\label{rem:split-bridge}
The accuracy decomposition (Eq.~\ref{eq:acc-decomp}) identifies two
levers for improving thinking-mode performance at budget $b$:
(i)~increase $F_L(b)$ (allow more chains to complete), or
(ii)~increase $\alpha_t(b)$ (extract more value from truncated chains).
Lever~(i) requires either larger budgets or shorter chains (model retraining).
Lever~(ii) can be achieved at inference time: by feeding a truncated
reasoning trace to a separate answer-generation pass, we replace
$\alpha_t(b)$ with $\alpha_{\mathrm{extract}}(b) \gg \alpha_t(b)$,
recovering value that would otherwise be lost to truncation waste.
\end{remark}

\noindent
This observation motivates \emph{split-budget generation}
(\S\ref{sec:method:split}): rather than attempting to prevent
truncation, we \emph{accept} it and recover its value through a
decoupled answer pass.
The truncation-waste framework provides the theoretical foundation
(identifying the coupling constraint as the sole mechanism behind the
tax), while the split-budget framework provides the operational remedy
(eliminating the constraint by separating reasoning and answering
into independent passes).

% ---------------------------------------------------------------------------
\subsection{The Coupling Tax as a Rate Constraint}
\label{sec:theory:coupling}
% ---------------------------------------------------------------------------

We now formalize the coupling constraint, yielding a tighter
characterization of \emph{why} split-budget generation helps and
\emph{how much} it can recover.
The following results are consequences of the decomposition framework
above; their value lies in making \emph{quantitative predictions} about
the gains from decoupling, which we verify empirically.

\begin{definition}[Coupled vs.\ split generation]
\label{def:coupled-split}
In \emph{coupled} generation, reasoning tokens~$Z$ and answer tokens~$A$
share a single budget: $|Z| + |A| \le b$.
In \emph{split-budget} generation, reasoning receives budget~$b_r$ and
answering receives an independent budget~$b_a$, with no constraint
coupling the two.
\end{definition}

\noindent
The coupling constraint $|Z|+|A| \le b$ acts as a \emph{rate limit}
on a joint channel: reasoning and answering compete for shared bandwidth.
When the reasoning trace is long ($|Z| \gg b - b_{\min}$, where $b_{\min}$
is the minimum tokens needed for a well-formed answer), the answer is
forcibly truncated or absent---not because the model lacks the answer, but
because the channel has no remaining capacity.
Split-budget generation removes this contention by allocating
independent channels.

\begin{proposition}[Recoverable coupling tax]
\label{prop:recoverable-tax}
Let $\alpha_{\mathrm{extract}}(b_r, b_a)$ denote the accuracy of a
separate answer-extraction pass that receives the first $b_r$ tokens
of the reasoning trace and generates up to $b_a$ answer tokens.
The accuracy gain from split-budget generation over coupled generation
at total budget $b$ satisfies:
\begin{equation}
  \Delta_{\mathrm{split}}(b)
  \;=\;
  \bigl(1 - F_L(b_r)\bigr) \cdot
  \bigl(\alpha_{\mathrm{extract}}(b_r, b_a) - \alpha_t(b)\bigr),
  \label{eq:recoverable-tax}
\end{equation}
where $b = b_r + b_a$.
This quantity is non-negative whenever
$\alpha_{\mathrm{extract}} \ge \alpha_t$---i.e., whenever the
decoupled answer pass extracts at least as much value as the truncated
coupled output---and grows with the truncation rate $(1 - F_L(b_r))$.
\end{proposition}

\begin{proof}
Split-budget accuracy decomposes as:
$\mathrm{Acc}_{\mathrm{split}}(b_r, b_a)
= F_L(b_r) \cdot \alpha_c
+ (1 - F_L(b_r)) \cdot \alpha_{\mathrm{extract}}(b_r, b_a)$.
Subtracting the coupled accuracy (Eq.~\ref{eq:acc-decomp}) at budget~$b$:
$\Delta_{\mathrm{split}}
= (1 - F_L(b_r))(\alpha_{\mathrm{extract}} - \alpha_t)
+ F_L(b_r) \cdot \alpha_c - F_L(b) \cdot \alpha_c$.
Since $b_r \le b$, we have $F_L(b_r) \le F_L(b)$; the second
difference is non-positive.
The first term dominates when truncation is high and
$\alpha_{\mathrm{extract}} \gg \alpha_t$, as in practice.
At budget $b_r = b$ (no separate answer budget), the first term
is exactly $(1 - F_L(b))(\alpha_{\mathrm{extract}} - \alpha_t)$.
\end{proof}

\noindent
\textbf{Empirical verification.}
On MATH-500 at $B_{\mathrm{think}}{=}2048$:
$1 - F_L(2048) = 0.912$ (91.2\% truncated), $\alpha_t \approx 0.365$
(from projection), and IRIS achieves 67.2\% vs.\ TOWN's 55.0\%
($\Delta = +12.2$\,pp, $p < 10^{-6}$).
At $B_{\mathrm{think}}{=}4096$: $1 - F_L(4096) = 0.553$, and the gap
narrows to $+2.2$\,pp ($p = 0.28$), \emph{exactly as predicted}:
lower truncation rate $\implies$ less recoverable waste
$\implies$ smaller $\Delta_{\mathrm{split}}$.
The ratio of truncation rates ($0.912/0.553 = 1.65$) directionally matches
the ratio of observed gaps, confirming that the coupling tax is the
\emph{sole mechanism} driving IRIS's advantage.

\begin{proposition}[Optimal budget allocation]
\label{prop:optimal-split}
Given a total token budget $B$ and the choice to allocate
$b_r$ tokens to reasoning and $b_a = B - b_r$ to answering,
the optimal reasoning budget is:
\begin{equation}
  b_r^* \;=\; \arg\max_{b_r \in [0, B]} \;
  F_L(b_r) \cdot \alpha_c
  + \bigl(1 - F_L(b_r)\bigr) \cdot
    \alpha_{\mathrm{extract}}(b_r, B - b_r).
  \label{eq:optimal-split}
\end{equation}
When $F_L$ is a log-normal CDF (as empirically observed;
see Table~\ref{tab:theory-verification}) and
$\alpha_{\mathrm{extract}}$ is monotonically non-decreasing in $b_r$
and concave in $b_a$, the objective has a unique interior maximum.
\end{proposition}

\noindent
The intuition is a bias--variance tradeoff: allocating more tokens to
reasoning increases the completion rate~$F_L(b_r)$ (reducing the
``truncation rate'' and its associated accuracy penalty), but leaves fewer
tokens for the answer pass when truncation does occur.
The optimum balances these competing demands.
In practice, our IRIS experiments use fixed $b_a = 256$ (sufficient for
answer extraction on MATH-500), allocating the remaining budget to reasoning.

\noindent
\textbf{Cross-architecture generality.}
The coupling constraint $|Z|+|A| \le b$ is not model-specific: it
applies to any autoregressive generator that produces reasoning and
answers in a single pass.
On DeepSeek-R1-Distill-Llama-8B (GSM8K, $n{=}40$), the same pattern
holds: at $b{=}256$, 82.5\% of chains are truncated vs.\ 7.5\% at
$b{=}512$; at $b{=}1024$, all chains complete naturally ($F_L = 1.0$,
$\alpha_c = 0.60$).
On MATH-500, DeepSeek-R1's natural-stop rate rises from 31.9\% at
$b{=}1024$ to 75.6\% at $b{=}4096$ (averaged over 4 seeds, $n{=}160$
total), mirroring the Qwen pattern.
The coupling tax is a \emph{structural property of the generation format},
not a model-family artifact.

% ===========================================================================

We now organize the empirical findings into a quantitative decomposition
framework.
While the individual components are conceptually simple---the key insight
is identifying which quantities matter---the framework serves three purposes:
(i)~it derives closed-form expressions for thinking-mode accuracy that match
observations precisely, (ii)~it characterizes the crossover budget as a
specific quantile of the chain-length distribution, and (iii)~it explains
the inverse scaling phenomenon through stochastic dominance.
We emphasize that our decomposition is an \emph{accounting framework} rather
than a deep theoretical result; its value lies in identifying the key parameters
($F_L$, $\alpha_c$, $\alpha_t$) that govern the phenomenon and showing they
suffice for quantitative prediction.

% ---------------------------------------------------------------------------
\subsection{A Formal Model of Truncation Waste}
\label{sec:theory:truncation}
% ---------------------------------------------------------------------------

\begin{definition}[Thinking chain length]
\label{def:chain-length}
For a model $\mathcal{M}$ and question $q$, let $L(q) \in \mathbb{N}$
denote the \emph{natural chain length}---the total number of output tokens
the model would generate in thinking mode if unconstrained
(i.e., $b \to \infty$), including both the reasoning trace and the
final answer.
Let $F_L(t) \triangleq \Pr(L \le t)$ denote the CDF of $L$ over the
question distribution $q \sim \mathcal{Q}$.
\end{definition}

\begin{definition}[Truncation rate]
\label{def:truncation}
Given token budget $b$, the \emph{truncation rate} is
$\rho(b) \triangleq 1 - F_L(b) = \Pr(L > b)$.
\end{definition}

\noindent
The key structural property of thinking mode is that the reasoning trace
and final answer share a single token budget.
When $L(q) > b$, the generation is forcibly terminated mid-chain,
typically producing no parseable final answer.
We capture this as:

\begin{assumption}[Binary outcome structure]
\label{asm:binary}
When the chain completes naturally ($L \le b$), the model produces a
well-formed answer with accuracy $\alpha_c \triangleq
\Pr(\text{correct} \mid L \le b) \in (0,1]$.
When truncated ($L > b$), the residual accuracy is
$\alpha_t \triangleq \Pr(\text{correct} \mid L > b) \ll \alpha_c$.
\end{assumption}

\noindent
This assumption is strongly supported empirically: at $b{=}512$
on GSM8K (full test set, $n{=}1{,}319$, HF engine), $\alpha_c = 99.0\%$ among natural-stop samples vs.\
$\alpha_t = 31.8\%$ among truncated samples (Finding~2, Table~\ref{tab:natural-stop-oracle}).
At lower budgets, $\alpha_t$ approaches zero.

\begin{proposition}[Accuracy decomposition]
\label{prop:acc-decomp}
Under Assumption~\ref{asm:binary}, thinking-mode accuracy decomposes as:
\begin{equation}
  \mathrm{Acc}_{\mathrm{think}}(b)
  = F_L(b) \cdot \alpha_c
  + \bigl(1 - F_L(b)\bigr) \cdot \alpha_t.
  \label{eq:acc-decomp}
\end{equation}
\end{proposition}

\begin{proof}
By the law of total probability, conditioning on $\{L \le b\}$:
$\mathrm{Acc}_{\mathrm{think}}(b)
= \Pr(\mathrm{correct} \mid L{\le}b)\,\Pr(L{\le}b)
+ \Pr(\mathrm{correct} \mid L{>}b)\,\Pr(L{>}b)
= \alpha_c\,F_L(b) + \alpha_t\,(1 - F_L(b))$.
\end{proof}

\begin{proposition}[The Thinking Tax]
\label{thm:thinking-tax}
Let $\mathrm{Acc}_{\mathrm{nt}}(b)$ denote non-thinking-mode accuracy
at budget $b$.
Setting $\alpha_t = 0$ as a first-order approximation, the thinking tax
satisfies:
\begin{equation}
  \underbrace{\mathrm{Acc}_{\mathrm{nt}}(b)
    - \mathrm{Acc}_{\mathrm{think}}(b)}_{\text{Thinking Tax}}
  \;\approx\;
  \mathrm{Acc}_{\mathrm{nt}}(b) - F_L(b) \cdot \alpha_c.
  \label{eq:tax}
\end{equation}
In particular, when $F_L(b) \ll 1$ (most chains exceed budget), the
tax approaches $\mathrm{Acc}_{\mathrm{nt}}(b)$ itself.
\end{proposition}

\begin{remark}[Budget-dependent $\alpha_t$]
\label{rem:alpha-t}
In practice, $\alpha_t$ varies with $b$: at $b{=}128$, nearly all
truncated samples lack any answer ($\alpha_t \approx 0$);
at $b{=}512$, answer-extraction heuristics and projection passes
recover partial answers ($\alpha_t \approx 31.8\%$).
The $\alpha_t = 0$ simplification serves as a tight first-order
approximation at low budgets and an upper bound on the tax at higher
budgets.
The full model (Eq.~\ref{eq:acc-decomp}) with budget-specific $\alpha_t$
values produces accurate predictions across all regimes.
\end{remark}

\noindent
\textbf{Empirical verification.}
At $b{=}512$, $F_L(512) = 0.374$ for Qwen3-8B and $\alpha_t = 31.8\%$, giving
$0.374 \times 99.0\% + 0.626 \times 31.8\% = 37.0\% + 19.9\% = 56.9\%$;
the observed value is 56.9\% ($n{=}1{,}319$), matching the prediction exactly.
At $b{=}128$, $F_L(128) \approx 0$ (0\% natural stop on the full test set), so
Eq.~\eqref{eq:acc-decomp} predicts $\mathrm{Acc}_{\mathrm{think}}(128) \approx \alpha_t \approx 3.0\%$;
the observed value is 3.0\%.

\noindent\textbf{Held-out prediction test.}
To test predictive power beyond in-sample verification, we estimate $\alpha_c$ and $\alpha_t$ from BBH at budget 512 ($\alpha_c{=}0.950$, $\alpha_t{=}0.275$) and use them to predict BBH accuracy at budgets 1024 and 2048, using only the observed completion rate $F_L$ at the target budget.
The framework predicts 69.3\% at $b{=}1024$ (observed: 73.6\%, error 4.3\,pp) and 86.8\% at $b{=}2048$ (observed: 86.0\%, error 0.8\,pp).
The larger error at $b{=}1024$ arises because $\alpha_t$ increases with budget (from 0.275 to 0.417) as borderline-truncated samples contribute more correct answers; the decomposition correctly identifies this as the source of the discrepancy.
This out-of-sample prediction---using parameters from one budget to predict accuracy at a different budget on the same benchmark---confirms the framework has genuine predictive utility beyond explaining known data.

\noindent
\textbf{Cross-benchmark prediction.}
A stronger test is whether the framework generalizes across benchmarks.
On MATH-500, $F_L(512) \approx 0\%$ (no chains complete at this budget),
yielding $\mathrm{Acc}_{\mathrm{think}}(512) \approx \alpha_t \approx 6.2\%$---consistent with the observed 6.2\%.
At $b{=}1024$, $F_L(1024) = 0.002$ (only 1/500 chains complete), with $\alpha_t = 0.178$ from partial answers recovered by projection, Eq.~\eqref{eq:acc-decomp} predicts $0.002 \times 1.0 + 0.998 \times 0.178 = 18.0\%$---matching the observed 18.0\% exactly.
At $b{=}2048$, $F_L(2048) = 0.178$ and $\alpha_c = 0.787$, $\alpha_t = 0.365$, predicting $0.178 \times 0.787 + 0.822 \times 0.365 = 44.0\%$---again an exact match (Table~\ref{tab:theory-verification}).
The framework correctly predicts that \emph{harder benchmarks amplify the tax}: longer expected chain lengths on MATH-500 push the crossover budget further right compared to GSM8K (${\sim}2048$), consistent with the prediction that $b^*$ scales with the chain-length CDF.
Table~\ref{tab:theory-verification} in the Appendix compares predicted
vs.\ observed accuracy across all configurations.

% ---------------------------------------------------------------------------
\subsection{Crossover Budget Characterization}
\label{sec:theory:crossover}
% ---------------------------------------------------------------------------

\begin{proposition}[Crossover budget]
\label{prop:crossover}
The crossover budget $b^*$---the smallest $b$ where thinking matches
non-thinking---satisfies:
\begin{equation}
  b^* \;\approx\; F_L^{-1}\!\!\left(
    \frac{\bar{\alpha}_{\mathrm{nt}}}{\alpha_c}\right),
  \label{eq:crossover}
\end{equation}
where $\bar{\alpha}_{\mathrm{nt}}$ is the saturated non-thinking accuracy
(for budgets beyond which $\mathrm{Acc}_{\mathrm{nt}}$ is approximately
constant).
That is, $b^*$ is the $(\bar{\alpha}_{\mathrm{nt}}/\alpha_c)$-quantile
of the chain-length distribution.
\end{proposition}

\begin{proof}
At $b^*$, setting $\alpha_t = 0$:
$\mathrm{Acc}_{\mathrm{think}}(b^*) = F_L(b^*) \cdot \alpha_c
= \mathrm{Acc}_{\mathrm{nt}}(b^*)$.
Since $\mathrm{Acc}_{\mathrm{nt}}$ saturates at $\bar{\alpha}_{\mathrm{nt}}$
for $b \ge b_{\mathrm{sat}}$ (empirically $\bar{\alpha}_{\mathrm{nt}} = 95.5\%$
at $b_{\mathrm{sat}} = 512$ for 27B; confirmed by 8B \texttt{nothink@1024}${=}$93.1\%, identical to \texttt{nothink@512}), we obtain
$F_L(b^*) = \bar{\alpha}_{\mathrm{nt}} / \alpha_c \approx 0.955/0.963 = 0.992$.
The crossover requires ${>}$99\% of chains to complete---i.e., $b^*$ is near the
99th percentile of $L$.
On GSM8K with 27B, whose chains are substantially longer, the crossover is expected well beyond 2048 tokens,
consistent with the observation that thinking has not matched non-thinking at any tested budget.
For 8B, with shorter chains, the crossover occurs at a lower budget but still exceeds 512 tokens.
\end{proof}

\begin{corollary}[Budget multiplier]
\label{cor:multiplier}
The budget multiplier $\gamma \triangleq b^*/b_{\mathrm{sat}}$ quantifies
the token overhead for thinking to become viable.
Since $b_{\mathrm{sat}}$ is determined by answer length (not chain length):
$\gamma = F_L^{-1}(\bar{\alpha}_{\mathrm{nt}}/\alpha_c) \big/ b_{\mathrm{sat}}$.
On MATH-500: $\gamma > 2048/1024 = 2\times$, since thinking has not yet caught non-thinking at $b{=}2048$ (Table~\ref{tab:math500-tax-full}), confirming that harder tasks demand larger multipliers.
For 27B on GSM8K, $b_{\mathrm{sat}} = 512$ (nothink saturates at 95.5\%) and the crossover is not reached at the tested budgets, implying $\gamma \gg 1$.
\end{corollary}

% ---------------------------------------------------------------------------
\subsection{Natural Stop as a Confidence Oracle}
\label{sec:theory:oracle}
% ---------------------------------------------------------------------------

\begin{proposition}[Oracle precision]
\label{prop:oracle}
The natural-stop event $\mathcal{S}(b) \triangleq \{|y| < b\}$
(generation terminates before exhausting budget) has positive predictive
value:
\begin{equation}
  \mathrm{PPV}(\mathcal{S}) \;=\;
  \Pr(\mathrm{correct} \mid \mathcal{S}(b))
  \;=\; \alpha_c.
  \label{eq:ppv}
\end{equation}
For non-thinking mode, $\mathrm{PPV}(\mathcal{S}_{\mathrm{nt}}) \ge \alpha_c$,
since early completion in non-thinking mode selects for easier questions
where model accuracy is higher.
\end{proposition}

\noindent
By Hoeffding's inequality, with $n{=}475$ natural-stop samples (8B, $b{=}512$)
and $\hat{\alpha}_c = 0.963$:
the 95\% lower confidence bound is
$\alpha_c \ge 0.963 - \sqrt{\ln(2/0.05)/(2 \cdot 475)} = 0.919$.
The signal is reliable even under finite-sample uncertainty.

\begin{remark}[Information-theoretic interpretation]
Early stopping implies the model's conditional entropy
$H(A \mid y_{1:t^*}) \approx 0$ at $t^* \ll b$---the model is
near-certain of its answer.
Unlike logit-based confidence~\citep{kadavath2022language}, this signal
is purely behavioral: it requires no access to model internals,
no calibration, and no additional compute.
\end{remark}

% ---------------------------------------------------------------------------
\subsection{Inverse Scaling with Model Size}
\label{sec:theory:scaling}
% ---------------------------------------------------------------------------

\begin{proposition}[Inverse scaling of the thinking tax]
\label{prop:inverse-scaling}
Let $L_M$ denote the chain-length distribution for a model of size $M$.
If $L_{M_2}$ stochastically dominates $L_{M_1}$ for $M_2 > M_1$
(i.e., $F_{L_{M_2}}(b) \le F_{L_{M_1}}(b)$ for all $b$), and
non-thinking accuracy is approximately size-invariant at moderate budgets,
then for any fixed budget $b$:
\begin{equation}
  \mathrm{Tax}(M_2, b) \;\ge\; \mathrm{Tax}(M_1, b).
  \label{eq:inverse-scaling}
\end{equation}
\end{proposition}

\begin{proof}
From Proposition~\ref{thm:thinking-tax}:
$\mathrm{Tax}(M, b) \approx \mathrm{Acc}_{\mathrm{nt}}(b) -
F_{L_M}(b) \cdot \alpha_c$.
With $\mathrm{Acc}_{\mathrm{nt}}$ approximately constant across $M$,
and $F_{L_{M_2}}(b) \le F_{L_{M_1}}(b)$, the second term decreases,
increasing the tax.
\end{proof}

\noindent
\textbf{Empirical verification.}
At $b{=}512$: 8B natural-stop rate = 37.4\%, 9B $\approx$ 0.5\%, 27B = 0.7\%.
Meanwhile, non-thinking accuracy at budget 512 is uniformly high: 8B = 93.1\%, 9B = 93.2\%, 27B = 95.5\%.
The thinking tax for 9B is $93.2\% - 15.5\% = 77.7$\,pp and for 27B is $95.5\% - 18.3\% = 77.2$\,pp, vs.\
8B's $93.1\% - 56.9\% = 36.2$\,pp---a ${\sim}2.1\times$ amplification from 8B to 9B,
as predicted by stochastic dominance of chain-length distributions.

\begin{corollary}[Crossover budget grows with model size]
\label{cor:crossover-scaling}
Since $b^*(M) \approx F_{L_M}^{-1}(\bar{\alpha}_{\mathrm{nt}}/\alpha_c)$
and $F_{L_M}^{-1}$ shifts right with $M$:
$M_2 > M_1 \implies b^*(M_2) \ge b^*(M_1)$.
Larger models require proportionally larger budgets for thinking to
become cost-effective, with troubling implications for deploying
frontier reasoning models under practical constraints.
\end{corollary}

% ---------------------------------------------------------------------------
\subsection{From Truncation Waste to Split-Budget Generation}
\label{sec:theory:bridge}
% ---------------------------------------------------------------------------

The decomposition framework reveals the \emph{structural cause} of
the thinking tax: the coupling constraint $|Z| + |A| \le b$ forces
reasoning and answering to compete for the same token budget.
This suggests a direct remedy: \emph{decouple} the two stages by
allocating separate budgets.

\begin{remark}[From diagnosis to intervention]
\label{rem:split-bridge}
The accuracy decomposition (Eq.~\ref{eq:acc-decomp}) identifies two
levers for improving thinking-mode performance at budget $b$:
(i)~increase $F_L(b)$ (allow more chains to complete), or
(ii)~increase $\alpha_t(b)$ (extract more value from truncated chains).
Lever~(i) requires either larger budgets or shorter chains (model retraining).
Lever~(ii) can be achieved at inference time: by feeding a truncated
reasoning trace to a separate answer-generation pass, we replace
$\alpha_t(b)$ with $\alpha_{\mathrm{extract}}(b) \gg \alpha_t(b)$,
recovering value that would otherwise be lost to truncation waste.
\end{remark}

\noindent
This observation motivates \emph{split-budget generation}
(\S\ref{sec:method:split}): rather than attempting to prevent
truncation, we \emph{accept} it and recover its value through a
decoupled answer pass.
The truncation-waste framework provides the theoretical foundation
(identifying the coupling constraint as the sole mechanism behind the
tax), while the split-budget framework provides the operational remedy
(eliminating the constraint by separating reasoning and answering
into independent passes).

% ---------------------------------------------------------------------------
\subsection{The Coupling Tax as a Rate Constraint}
\label{sec:theory:coupling}
% ---------------------------------------------------------------------------

We now formalize the coupling constraint, yielding a tighter
characterization of \emph{why} split-budget generation helps and
\emph{how much} it can recover.
The following results are consequences of the decomposition framework
above; their value lies in making \emph{quantitative predictions} about
the gains from decoupling, which we verify empirically.

\begin{definition}[Coupled vs.\ split generation]
\label{def:coupled-split}
In \emph{coupled} generation, reasoning tokens~$Z$ and answer tokens~$A$
share a single budget: $|Z| + |A| \le b$.
In \emph{split-budget} generation, reasoning receives budget~$b_r$ and
answering receives an independent budget~$b_a$, with no constraint
coupling the two.
\end{definition}

\noindent
The coupling constraint $|Z|+|A| \le b$ acts as a \emph{rate limit}
on a joint channel: reasoning and answering compete for shared bandwidth.
When the reasoning trace is long ($|Z| \gg b - b_{\min}$, where $b_{\min}$
is the minimum tokens needed for a well-formed answer), the answer is
forcibly truncated or absent---not because the model lacks the answer, but
because the channel has no remaining capacity.
Split-budget generation removes this contention by allocating
independent channels.

\begin{proposition}[Recoverable coupling tax]
\label{prop:recoverable-tax}
Let $\alpha_{\mathrm{extract}}(b_r, b_a)$ denote the accuracy of a
separate answer-extraction pass that receives the first $b_r$ tokens
of the reasoning trace and generates up to $b_a$ answer tokens.
The accuracy gain from split-budget generation over coupled generation
at total budget $b$ satisfies:
\begin{equation}
  \Delta_{\mathrm{split}}(b)
  \;=\;
  \bigl(1 - F_L(b_r)\bigr) \cdot
  \bigl(\alpha_{\mathrm{extract}}(b_r, b_a) - \alpha_t(b)\bigr),
  \label{eq:recoverable-tax}
\end{equation}
where $b = b_r + b_a$.
This quantity is non-negative whenever
$\alpha_{\mathrm{extract}} \ge \alpha_t$---i.e., whenever the
decoupled answer pass extracts at least as much value as the truncated
coupled output---and grows with the truncation rate $(1 - F_L(b_r))$.
\end{proposition}

\begin{proof}
Split-budget accuracy decomposes as:
$\mathrm{Acc}_{\mathrm{split}}(b_r, b_a)
= F_L(b_r) \cdot \alpha_c
+ (1 - F_L(b_r)) \cdot \alpha_{\mathrm{extract}}(b_r, b_a)$.
Subtracting the coupled accuracy (Eq.~\ref{eq:acc-decomp}) at budget~$b$:
$\Delta_{\mathrm{split}}
= (1 - F_L(b_r))(\alpha_{\mathrm{extract}} - \alpha_t)
+ F_L(b_r) \cdot \alpha_c - F_L(b) \cdot \alpha_c$.
Since $b_r \le b$, we have $F_L(b_r) \le F_L(b)$; the second
difference is non-positive.
The first term dominates when truncation is high and
$\alpha_{\mathrm{extract}} \gg \alpha_t$, as in practice.
At budget $b_r = b$ (no separate answer budget), the first term
is exactly $(1 - F_L(b))(\alpha_{\mathrm{extract}} - \alpha_t)$.
\end{proof}

\noindent
\textbf{Empirical verification.}
On MATH-500 at $B_{\mathrm{think}}{=}2048$:
$1 - F_L(2048) = 0.912$ (91.2\% truncated), $\alpha_t \approx 0.365$
(from projection), and IRIS achieves 67.2\% vs.\ TOWN's 55.0\%
($\Delta = +12.2$\,pp, $p < 10^{-6}$).
At $B_{\mathrm{think}}{=}4096$: $1 - F_L(4096) = 0.553$, and the gap
narrows to $+2.2$\,pp ($p = 0.28$), \emph{exactly as predicted}:
lower truncation rate $\implies$ less recoverable waste
$\implies$ smaller $\Delta_{\mathrm{split}}$.
The ratio of truncation rates ($0.912/0.553 = 1.65$) directionally matches
the ratio of observed gaps, confirming that the coupling tax is the
\emph{sole mechanism} driving IRIS's advantage.

\begin{proposition}[Optimal budget allocation]
\label{prop:optimal-split}
Given a total token budget $B$ and the choice to allocate
$b_r$ tokens to reasoning and $b_a = B - b_r$ to answering,
the optimal reasoning budget is:
\begin{equation}
  b_r^* \;=\; \arg\max_{b_r \in [0, B]} \;
  F_L(b_r) \cdot \alpha_c
  + \bigl(1 - F_L(b_r)\bigr) \cdot
    \alpha_{\mathrm{extract}}(b_r, B - b_r).
  \label{eq:optimal-split}
\end{equation}
When $F_L$ is a log-normal CDF (as empirically observed;
see Table~\ref{tab:theory-verification}) and
$\alpha_{\mathrm{extract}}$ is monotonically non-decreasing in $b_r$
and concave in $b_a$, the objective has a unique interior maximum.
\end{proposition}

\noindent
The intuition is a bias--variance tradeoff: allocating more tokens to
reasoning increases the completion rate~$F_L(b_r)$ (reducing the
``truncation rate'' and its associated accuracy penalty), but leaves fewer
tokens for the answer pass when truncation does occur.
The optimum balances these competing demands.
In practice, our IRIS experiments use fixed $b_a = 256$ (sufficient for
answer extraction on MATH-500), allocating the remaining budget to reasoning.

\noindent
\textbf{Cross-architecture generality.}
The coupling constraint $|Z|+|A| \le b$ is not model-specific: it
applies to any autoregressive generator that produces reasoning and
answers in a single pass.
On DeepSeek-R1-Distill-Llama-8B (GSM8K, $n{=}40$), the same pattern
holds: at $b{=}256$, 82.5\% of chains are truncated vs.\ 7.5\% at
$b{=}512$; at $b{=}1024$, all chains complete naturally ($F_L = 1.0$,
$\alpha_c = 0.60$).
On MATH-500, DeepSeek-R1's natural-stop rate rises from 31.9\% at
$b{=}1024$ to 75.6\% at $b{=}4096$ (averaged over 4 seeds, $n{=}160$
total), mirroring the Qwen pattern.
The coupling tax is a \emph{structural property of the generation format},
not a model-family artifact.

% ===========================================================================

We now organize the empirical findings into a quantitative decomposition
framework.
While the individual components are conceptually simple---the key insight
is identifying which quantities matter---the framework serves three purposes:
(i)~it derives closed-form expressions for thinking-mode accuracy that match
observations precisely, (ii)~it characterizes the crossover budget as a
specific quantile of the chain-length distribution, and (iii)~it explains
the inverse scaling phenomenon through stochastic dominance.
We emphasize that our decomposition is an \emph{accounting framework} rather
than a deep theoretical result; its value lies in identifying the key parameters
($F_L$, $\alpha_c$, $\alpha_t$) that govern the phenomenon and showing they
suffice for quantitative prediction.

% ---------------------------------------------------------------------------
\subsection{A Formal Model of Truncation Waste}
\label{sec:theory:truncation}
% ---------------------------------------------------------------------------

\begin{definition}[Thinking chain length]
\label{def:chain-length}
For a model $\mathcal{M}$ and question $q$, let $L(q) \in \mathbb{N}$
denote the \emph{natural chain length}---the total number of output tokens
the model would generate in thinking mode if unconstrained
(i.e., $b \to \infty$), including both the reasoning trace and the
final answer.
Let $F_L(t) \triangleq \Pr(L \le t)$ denote the CDF of $L$ over the
question distribution $q \sim \mathcal{Q}$.
\end{definition}

\begin{definition}[Truncation rate]
\label{def:truncation}
Given token budget $b$, the \emph{truncation rate} is
$\rho(b) \triangleq 1 - F_L(b) = \Pr(L > b)$.
\end{definition}

\noindent
The key structural property of thinking mode is that the reasoning trace
and final answer share a single token budget.
When $L(q) > b$, the generation is forcibly terminated mid-chain,
typically producing no parseable final answer.
We capture this as:

\begin{assumption}[Binary outcome structure]
\label{asm:binary}
When the chain completes naturally ($L \le b$), the model produces a
well-formed answer with accuracy $\alpha_c \triangleq
\Pr(\text{correct} \mid L \le b) \in (0,1]$.
When truncated ($L > b$), the residual accuracy is
$\alpha_t \triangleq \Pr(\text{correct} \mid L > b) \ll \alpha_c$.
\end{assumption}

\noindent
This assumption is strongly supported empirically: at $b{=}512$
on GSM8K (full test set, $n{=}1{,}319$, HF engine), $\alpha_c = 99.0\%$ among natural-stop samples vs.\
$\alpha_t = 31.8\%$ among truncated samples (Finding~2, Table~\ref{tab:natural-stop-oracle}).
At lower budgets, $\alpha_t$ approaches zero.

\begin{proposition}[Accuracy decomposition]
\label{prop:acc-decomp}
Under Assumption~\ref{asm:binary}, thinking-mode accuracy decomposes as:
\begin{equation}
  \mathrm{Acc}_{\mathrm{think}}(b)
  = F_L(b) \cdot \alpha_c
  + \bigl(1 - F_L(b)\bigr) \cdot \alpha_t.
  \label{eq:acc-decomp}
\end{equation}
\end{proposition}

\begin{proof}
By the law of total probability, conditioning on $\{L \le b\}$:
$\mathrm{Acc}_{\mathrm{think}}(b)
= \Pr(\mathrm{correct} \mid L{\le}b)\,\Pr(L{\le}b)
+ \Pr(\mathrm{correct} \mid L{>}b)\,\Pr(L{>}b)
= \alpha_c\,F_L(b) + \alpha_t\,(1 - F_L(b))$.
\end{proof}

\begin{proposition}[The Thinking Tax]
\label{thm:thinking-tax}
Let $\mathrm{Acc}_{\mathrm{nt}}(b)$ denote non-thinking-mode accuracy
at budget $b$.
Setting $\alpha_t = 0$ as a first-order approximation, the thinking tax
satisfies:
\begin{equation}
  \underbrace{\mathrm{Acc}_{\mathrm{nt}}(b)
    - \mathrm{Acc}_{\mathrm{think}}(b)}_{\text{Thinking Tax}}
  \;\approx\;
  \mathrm{Acc}_{\mathrm{nt}}(b) - F_L(b) \cdot \alpha_c.
  \label{eq:tax}
\end{equation}
In particular, when $F_L(b) \ll 1$ (most chains exceed budget), the
tax approaches $\mathrm{Acc}_{\mathrm{nt}}(b)$ itself.
\end{proposition}

\begin{remark}[Budget-dependent $\alpha_t$]
\label{rem:alpha-t}
In practice, $\alpha_t$ varies with $b$: at $b{=}128$, nearly all
truncated samples lack any answer ($\alpha_t \approx 0$);
at $b{=}512$, answer-extraction heuristics and projection passes
recover partial answers ($\alpha_t \approx 31.8\%$).
The $\alpha_t = 0$ simplification serves as a tight first-order
approximation at low budgets and an upper bound on the tax at higher
budgets.
The full model (Eq.~\ref{eq:acc-decomp}) with budget-specific $\alpha_t$
values produces accurate predictions across all regimes.
\end{remark}

\noindent
\textbf{Empirical verification.}
At $b{=}512$, $F_L(512) = 0.374$ for Qwen3-8B and $\alpha_t = 31.8\%$, giving
$0.374 \times 99.0\% + 0.626 \times 31.8\% = 37.0\% + 19.9\% = 56.9\%$;
the observed value is 56.9\% ($n{=}1{,}319$), matching the prediction exactly.
At $b{=}128$, $F_L(128) \approx 0$ (0\% natural stop on the full test set), so
Eq.~\eqref{eq:acc-decomp} predicts $\mathrm{Acc}_{\mathrm{think}}(128) \approx \alpha_t \approx 3.0\%$;
the observed value is 3.0\%.

\noindent\textbf{Held-out prediction test.}
To test predictive power beyond in-sample verification, we estimate $\alpha_c$ and $\alpha_t$ from BBH at budget 512 ($\alpha_c{=}0.950$, $\alpha_t{=}0.275$) and use them to predict BBH accuracy at budgets 1024 and 2048, using only the observed completion rate $F_L$ at the target budget.
The framework predicts 69.3\% at $b{=}1024$ (observed: 73.6\%, error 4.3\,pp) and 86.8\% at $b{=}2048$ (observed: 86.0\%, error 0.8\,pp).
The larger error at $b{=}1024$ arises because $\alpha_t$ increases with budget (from 0.275 to 0.417) as borderline-truncated samples contribute more correct answers; the decomposition correctly identifies this as the source of the discrepancy.
This out-of-sample prediction---using parameters from one budget to predict accuracy at a different budget on the same benchmark---confirms the framework has genuine predictive utility beyond explaining known data.

\noindent
\textbf{Cross-benchmark prediction.}
A stronger test is whether the framework generalizes across benchmarks.
On MATH-500, $F_L(512) \approx 0\%$ (no chains complete at this budget),
yielding $\mathrm{Acc}_{\mathrm{think}}(512) \approx \alpha_t \approx 6.2\%$---consistent with the observed 6.2\%.
At $b{=}1024$, $F_L(1024) = 0.002$ (only 1/500 chains complete), with $\alpha_t = 0.178$ from partial answers recovered by projection, Eq.~\eqref{eq:acc-decomp} predicts $0.002 \times 1.0 + 0.998 \times 0.178 = 18.0\%$---matching the observed 18.0\% exactly.
At $b{=}2048$, $F_L(2048) = 0.178$ and $\alpha_c = 0.787$, $\alpha_t = 0.365$, predicting $0.178 \times 0.787 + 0.822 \times 0.365 = 44.0\%$---again an exact match (Table~\ref{tab:theory-verification}).
The framework correctly predicts that \emph{harder benchmarks amplify the tax}: longer expected chain lengths on MATH-500 push the crossover budget further right compared to GSM8K (${\sim}2048$), consistent with the prediction that $b^*$ scales with the chain-length CDF.
Table~\ref{tab:theory-verification} in the Appendix compares predicted
vs.\ observed accuracy across all configurations.

% ---------------------------------------------------------------------------
\subsection{Crossover Budget Characterization}
\label{sec:theory:crossover}
% ---------------------------------------------------------------------------

\begin{proposition}[Crossover budget]
\label{prop:crossover}
The crossover budget $b^*$---the smallest $b$ where thinking matches
non-thinking---satisfies:
\begin{equation}
  b^* \;\approx\; F_L^{-1}\!\!\left(
    \frac{\bar{\alpha}_{\mathrm{nt}}}{\alpha_c}\right),
  \label{eq:crossover}
\end{equation}
where $\bar{\alpha}_{\mathrm{nt}}$ is the saturated non-thinking accuracy
(for budgets beyond which $\mathrm{Acc}_{\mathrm{nt}}$ is approximately
constant).
That is, $b^*$ is the $(\bar{\alpha}_{\mathrm{nt}}/\alpha_c)$-quantile
of the chain-length distribution.
\end{proposition}

\begin{proof}
At $b^*$, setting $\alpha_t = 0$:
$\mathrm{Acc}_{\mathrm{think}}(b^*) = F_L(b^*) \cdot \alpha_c
= \mathrm{Acc}_{\mathrm{nt}}(b^*)$.
Since $\mathrm{Acc}_{\mathrm{nt}}$ saturates at $\bar{\alpha}_{\mathrm{nt}}$
for $b \ge b_{\mathrm{sat}}$ (empirically $\bar{\alpha}_{\mathrm{nt}} = 95.5\%$
at $b_{\mathrm{sat}} = 512$ for 27B; confirmed by 8B \texttt{nothink@1024}${=}$93.1\%, identical to \texttt{nothink@512}), we obtain
$F_L(b^*) = \bar{\alpha}_{\mathrm{nt}} / \alpha_c \approx 0.955/0.963 = 0.992$.
The crossover requires ${>}$99\% of chains to complete---i.e., $b^*$ is near the
99th percentile of $L$.
On GSM8K with 27B, whose chains are substantially longer, the crossover is expected well beyond 2048 tokens,
consistent with the observation that thinking has not matched non-thinking at any tested budget.
For 8B, with shorter chains, the crossover occurs at a lower budget but still exceeds 512 tokens.
\end{proof}

\begin{corollary}[Budget multiplier]
\label{cor:multiplier}
The budget multiplier $\gamma \triangleq b^*/b_{\mathrm{sat}}$ quantifies
the token overhead for thinking to become viable.
Since $b_{\mathrm{sat}}$ is determined by answer length (not chain length):
$\gamma = F_L^{-1}(\bar{\alpha}_{\mathrm{nt}}/\alpha_c) \big/ b_{\mathrm{sat}}$.
On MATH-500: $\gamma > 2048/1024 = 2\times$, since thinking has not yet caught non-thinking at $b{=}2048$ (Table~\ref{tab:math500-tax-full}), confirming that harder tasks demand larger multipliers.
For 27B on GSM8K, $b_{\mathrm{sat}} = 512$ (nothink saturates at 95.5\%) and the crossover is not reached at the tested budgets, implying $\gamma \gg 1$.
\end{corollary}

% ---------------------------------------------------------------------------
\subsection{Natural Stop as a Confidence Oracle}
\label{sec:theory:oracle}
% ---------------------------------------------------------------------------

\begin{proposition}[Oracle precision]
\label{prop:oracle}
The natural-stop event $\mathcal{S}(b) \triangleq \{|y| < b\}$
(generation terminates before exhausting budget) has positive predictive
value:
\begin{equation}
  \mathrm{PPV}(\mathcal{S}) \;=\;
  \Pr(\mathrm{correct} \mid \mathcal{S}(b))
  \;=\; \alpha_c.
  \label{eq:ppv}
\end{equation}
For non-thinking mode, $\mathrm{PPV}(\mathcal{S}_{\mathrm{nt}}) \ge \alpha_c$,
since early completion in non-thinking mode selects for easier questions
where model accuracy is higher.
\end{proposition}

\noindent
By Hoeffding's inequality, with $n{=}475$ natural-stop samples (8B, $b{=}512$)
and $\hat{\alpha}_c = 0.963$:
the 95\% lower confidence bound is
$\alpha_c \ge 0.963 - \sqrt{\ln(2/0.05)/(2 \cdot 475)} = 0.919$.
The signal is reliable even under finite-sample uncertainty.

\begin{remark}[Information-theoretic interpretation]
Early stopping implies the model's conditional entropy
$H(A \mid y_{1:t^*}) \approx 0$ at $t^* \ll b$---the model is
near-certain of its answer.
Unlike logit-based confidence~\citep{kadavath2022language}, this signal
is purely behavioral: it requires no access to model internals,
no calibration, and no additional compute.
\end{remark}

% ---------------------------------------------------------------------------
\subsection{Inverse Scaling with Model Size}
\label{sec:theory:scaling}
% ---------------------------------------------------------------------------

\begin{proposition}[Inverse scaling of the thinking tax]
\label{prop:inverse-scaling}
Let $L_M$ denote the chain-length distribution for a model of size $M$.
If $L_{M_2}$ stochastically dominates $L_{M_1}$ for $M_2 > M_1$
(i.e., $F_{L_{M_2}}(b) \le F_{L_{M_1}}(b)$ for all $b$), and
non-thinking accuracy is approximately size-invariant at moderate budgets,
then for any fixed budget $b$:
\begin{equation}
  \mathrm{Tax}(M_2, b) \;\ge\; \mathrm{Tax}(M_1, b).
  \label{eq:inverse-scaling}
\end{equation}
\end{proposition}

\begin{proof}
From Proposition~\ref{thm:thinking-tax}:
$\mathrm{Tax}(M, b) \approx \mathrm{Acc}_{\mathrm{nt}}(b) -
F_{L_M}(b) \cdot \alpha_c$.
With $\mathrm{Acc}_{\mathrm{nt}}$ approximately constant across $M$,
and $F_{L_{M_2}}(b) \le F_{L_{M_1}}(b)$, the second term decreases,
increasing the tax.
\end{proof}

\noindent
\textbf{Empirical verification.}
At $b{=}512$: 8B natural-stop rate = 37.4\%, 9B $\approx$ 0.5\%, 27B = 0.7\%.
Meanwhile, non-thinking accuracy at budget 512 is uniformly high: 8B = 93.1\%, 9B = 93.2\%, 27B = 95.5\%.
The thinking tax for 9B is $93.2\% - 15.5\% = 77.7$\,pp and for 27B is $95.5\% - 18.3\% = 77.2$\,pp, vs.\
8B's $93.1\% - 56.9\% = 36.2$\,pp---a ${\sim}2.1\times$ amplification from 8B to 9B,
as predicted by stochastic dominance of chain-length distributions.

\begin{corollary}[Crossover budget grows with model size]
\label{cor:crossover-scaling}
Since $b^*(M) \approx F_{L_M}^{-1}(\bar{\alpha}_{\mathrm{nt}}/\alpha_c)$
and $F_{L_M}^{-1}$ shifts right with $M$:
$M_2 > M_1 \implies b^*(M_2) \ge b^*(M_1)$.
Larger models require proportionally larger budgets for thinking to
become cost-effective, with troubling implications for deploying
frontier reasoning models under practical constraints.
\end{corollary}

% ---------------------------------------------------------------------------
\subsection{From Truncation Waste to Split-Budget Generation}
\label{sec:theory:bridge}
% ---------------------------------------------------------------------------

The decomposition framework reveals the \emph{structural cause} of
the thinking tax: the coupling constraint $|Z| + |A| \le b$ forces
reasoning and answering to compete for the same token budget.
This suggests a direct remedy: \emph{decouple} the two stages by
allocating separate budgets.

\begin{remark}[From diagnosis to intervention]
\label{rem:split-bridge}
The accuracy decomposition (Eq.~\ref{eq:acc-decomp}) identifies two
levers for improving thinking-mode performance at budget $b$:
(i)~increase $F_L(b)$ (allow more chains to complete), or
(ii)~increase $\alpha_t(b)$ (extract more value from truncated chains).
Lever~(i) requires either larger budgets or shorter chains (model retraining).
Lever~(ii) can be achieved at inference time: by feeding a truncated
reasoning trace to a separate answer-generation pass, we replace
$\alpha_t(b)$ with $\alpha_{\mathrm{extract}}(b) \gg \alpha_t(b)$,
recovering value that would otherwise be lost to truncation waste.
\end{remark}

\noindent
This observation motivates \emph{split-budget generation}
(\S\ref{sec:method:split}): rather than attempting to prevent
truncation, we \emph{accept} it and recover its value through a
decoupled answer pass.
The truncation-waste framework provides the theoretical foundation
(identifying the coupling constraint as the sole mechanism behind the
tax), while the split-budget framework provides the operational remedy
(eliminating the constraint by separating reasoning and answering
into independent passes).

% ---------------------------------------------------------------------------
\subsection{The Coupling Tax as a Rate Constraint}
\label{sec:theory:coupling}
% ---------------------------------------------------------------------------

We now formalize the coupling constraint, yielding a tighter
characterization of \emph{why} split-budget generation helps and
\emph{how much} it can recover.
The following results are consequences of the decomposition framework
above; their value lies in making \emph{quantitative predictions} about
the gains from decoupling, which we verify empirically.

\begin{definition}[Coupled vs.\ split generation]
\label{def:coupled-split}
In \emph{coupled} generation, reasoning tokens~$Z$ and answer tokens~$A$
share a single budget: $|Z| + |A| \le b$.
In \emph{split-budget} generation, reasoning receives budget~$b_r$ and
answering receives an independent budget~$b_a$, with no constraint
coupling the two.
\end{definition}

\noindent
The coupling constraint $|Z|+|A| \le b$ acts as a \emph{rate limit}
on a joint channel: reasoning and answering compete for shared bandwidth.
When the reasoning trace is long ($|Z| \gg b - b_{\min}$, where $b_{\min}$
is the minimum tokens needed for a well-formed answer), the answer is
forcibly truncated or absent---not because the model lacks the answer, but
because the channel has no remaining capacity.
Split-budget generation removes this contention by allocating
independent channels.

\begin{proposition}[Recoverable coupling tax]
\label{prop:recoverable-tax}
Let $\alpha_{\mathrm{extract}}(b_r, b_a)$ denote the accuracy of a
separate answer-extraction pass that receives the first $b_r$ tokens
of the reasoning trace and generates up to $b_a$ answer tokens.
The accuracy gain from split-budget generation over coupled generation
at total budget $b$ satisfies:
\begin{equation}
  \Delta_{\mathrm{split}}(b)
  \;=\;
  \bigl(1 - F_L(b_r)\bigr) \cdot
  \bigl(\alpha_{\mathrm{extract}}(b_r, b_a) - \alpha_t(b)\bigr),
  \label{eq:recoverable-tax}
\end{equation}
where $b = b_r + b_a$.
This quantity is non-negative whenever
$\alpha_{\mathrm{extract}} \ge \alpha_t$---i.e., whenever the
decoupled answer pass extracts at least as much value as the truncated
coupled output---and grows with the truncation rate $(1 - F_L(b_r))$.
\end{proposition}

\begin{proof}
Split-budget accuracy decomposes as:
$\mathrm{Acc}_{\mathrm{split}}(b_r, b_a)
= F_L(b_r) \cdot \alpha_c
+ (1 - F_L(b_r)) \cdot \alpha_{\mathrm{extract}}(b_r, b_a)$.
Subtracting the coupled accuracy (Eq.~\ref{eq:acc-decomp}) at budget~$b$:
$\Delta_{\mathrm{split}}
= (1 - F_L(b_r))(\alpha_{\mathrm{extract}} - \alpha_t)
+ F_L(b_r) \cdot \alpha_c - F_L(b) \cdot \alpha_c$.
Since $b_r \le b$, we have $F_L(b_r) \le F_L(b)$; the second
difference is non-positive.
The first term dominates when truncation is high and
$\alpha_{\mathrm{extract}} \gg \alpha_t$, as in practice.
At budget $b_r = b$ (no separate answer budget), the first term
is exactly $(1 - F_L(b))(\alpha_{\mathrm{extract}} - \alpha_t)$.
\end{proof}

\noindent
\textbf{Empirical verification.}
On MATH-500 at $B_{\mathrm{think}}{=}2048$:
$1 - F_L(2048) = 0.912$ (91.2\% truncated), $\alpha_t \approx 0.365$
(from projection), and IRIS achieves 67.2\% vs.\ TOWN's 55.0\%
($\Delta = +12.2$\,pp, $p < 10^{-6}$).
At $B_{\mathrm{think}}{=}4096$: $1 - F_L(4096) = 0.553$, and the gap
narrows to $+2.2$\,pp ($p = 0.28$), \emph{exactly as predicted}:
lower truncation rate $\implies$ less recoverable waste
$\implies$ smaller $\Delta_{\mathrm{split}}$.
The ratio of truncation rates ($0.912/0.553 = 1.65$) directionally matches
the ratio of observed gaps, confirming that the coupling tax is the
\emph{sole mechanism} driving IRIS's advantage.

\begin{proposition}[Optimal budget allocation]
\label{prop:optimal-split}
Given a total token budget $B$ and the choice to allocate
$b_r$ tokens to reasoning and $b_a = B - b_r$ to answering,
the optimal reasoning budget is:
\begin{equation}
  b_r^* \;=\; \arg\max_{b_r \in [0, B]} \;
  F_L(b_r) \cdot \alpha_c
  + \bigl(1 - F_L(b_r)\bigr) \cdot
    \alpha_{\mathrm{extract}}(b_r, B - b_r).
  \label{eq:optimal-split}
\end{equation}
When $F_L$ is a log-normal CDF (as empirically observed;
see Table~\ref{tab:theory-verification}) and
$\alpha_{\mathrm{extract}}$ is monotonically non-decreasing in $b_r$
and concave in $b_a$, the objective has a unique interior maximum.
\end{proposition}

\noindent
The intuition is a bias--variance tradeoff: allocating more tokens to
reasoning increases the completion rate~$F_L(b_r)$ (reducing the
``truncation rate'' and its associated accuracy penalty), but leaves fewer
tokens for the answer pass when truncation does occur.
The optimum balances these competing demands.
In practice, our IRIS experiments use fixed $b_a = 256$ (sufficient for
answer extraction on MATH-500), allocating the remaining budget to reasoning.

\noindent
\textbf{Cross-architecture generality.}
The coupling constraint $|Z|+|A| \le b$ is not model-specific: it
applies to any autoregressive generator that produces reasoning and
answers in a single pass.
On DeepSeek-R1-Distill-Llama-8B (GSM8K, $n{=}40$), the same pattern
holds: at $b{=}256$, 82.5\% of chains are truncated vs.\ 7.5\% at
$b{=}512$; at $b{=}1024$, all chains complete naturally ($F_L = 1.0$,
$\alpha_c = 0.60$).
On MATH-500, DeepSeek-R1's natural-stop rate rises from 31.9\% at
$b{=}1024$ to 75.6\% at $b{=}4096$ (averaged over 4 seeds, $n{=}160$
total), mirroring the Qwen pattern.
The coupling tax is a \emph{structural property of the generation format},
not a model-family artifact.
